
\section{Relaciones con los módulos libres}
\label{IV.4}

\begin{proposition}
\label{IV.4.1}
Sean~$A$ un anillo, $I$ un ideal de~$A$ y~$M$ un $A$-módulo.
Supongamos que se cumple una de las dos hipótesis siguientes:
\begin{enumerate}
\item[(a)] $I$ es nilpotente.
\item[(b)] $A$ es noetheriano, $I$ está contenido en el radical de
Jacobson de~$A$ y~$M$ es de tipo finito.
\end{enumerate}
Para que~$M$ sea libre sobre~$A$, es necesario y suficiente que
$M\otimes_A A/I$ sea libre sobre~$A/I$ y que
$\Tor_1^A(M,A/I)=0$.
\end{proposition}

La necesidad es clara;
\marginpar{94}
demostremos la suficiencia. Sea~$(e_i)$ una familia
de elementos de~$M$ cuya imagen en~$M\otimes_A A/I=M/IM$ define una base
sobre~$A/I$ (la familia es finita en el caso~(b)). Sea~$L$ el $A$-módulo
libre construido sobre el mismo conjunto de índices. Se tiene entonces un
homomorfismo~$L\to M$ tal que el funtor~$T$ de tensorización por~$A/I$
induce un isomorfismo~$T(L)\isomto T(M)$. Si~$Q$ es el conúcleo de
$L\to M$, se tiene~$T(Q)=0$, de donde~$Q=0$ por Nakayama (válido bajo
cualquiera de las condiciones~(a) y~(b)). Por tanto, $L\to M$ es
sobreyectivo. Sea~$R$ su núcleo; se tiene una sucesión exacta
$$
0\to R\to L\to M\to0.
$$
Como~$\Tor_1^A(M,A/I)=0$, se obtiene una sucesión exacta
$0\to T(R)\to T(L)\to T(M)\to0$, de donde~$T(R)=0$ y, de nuevo,
\ifthenelse{\boolean{orig}}{$T(R)=0$}{$R=0$}
por Nakayama (teniendo en cuenta que, en el caso~(b), $R$ es de tipo
finito porque~$A$ se supuso noetheriano).

\ifthenelse{\boolean{orig}}
{\begin{corollairestar}}
{\begin{corollaire}\label{IV.4.2}}
La condición~$\Tor_1^A(M,A/I)=0$ puede sustituirse por la siguiente: el
homomorfismo canónico sobreyectivo
\begin{equation*}
\ifthenelse{\boolean{orig}}
{\gdef\IVFourStarLabel{eq:IV.2.*}}
{\gdef\IVFourStarLabel{eq:IV.4.*}}
\label{\IVFourStarLabel}
\tag{$*$} {\gr_I^0(M)\otimes_{A/I}\gr_I(A)\to\gr_I(M)}
\end{equation*}
es un isomorfismo.
\ifthenelse{\boolean{orig}}
{\end{corollairestar}}
{\end{corollaire}}

En efecto, si~$M$ es libre, esto se cumple sin duda. Hay que demostrar,
pues, que si~$M\otimes_A A/I$ es libre sobre~$A/I$ y se cumple la condición
sobre los graduados asociados, entonces~$M$ es libre. Se retoma la
demostración anterior construyendo~$L\to M$. La hipótesis implica que este
homomorfismo induce un isomorfismo entre los graduados asociados; por
tanto, su núcleo está contenido en la intersección de los~$I^nL$ y es nulo
(como es trivial en~(a) y bien conocido en~(b)). c.q.d.

\ifthenelse{\boolean{orig}}{\refstepcounter{subsection}}{}

\begin{corollaire}
\label{IV.4.3}
Supongamos que~$A/I$ es un cuerpo. Entonces son equivalentes las condiciones
siguientes sobre~$M$:
\begin{enumerate}
\item[(i)] $M$ es libre.
\item[(ii)] $M$ es proyectivo.
\item[(iii)] $M$ es plano.
\item[(iv)] $\Tor_1^A(M,A/I)=0$.
\item[(v)] El homomorfismo canónico~\eqref{\IVFourStarLabel}
es biyectivo.
\end{enumerate}
\end{corollaire}

En efecto, en el caso considerado, $M\otimes_A A/I$ es automáticamente
libre.

El resultado anterior es válido en los dos casos siguientes:
\begin{enumerate}
\item[(a)] $M$ es un módulo \emph{arbitrario} sobre un anillo local~$A$
cuyo ideal maximal~$I$ es \emph{nilpotente} (por ejemplo, un anillo local
artiniano).
\item[(b)] $M$ es un módulo \emph{de tipo finito} sobre un anillo
\emph{local noetheriano}.
\end{enumerate}

Recordemos, a modo de referencia:
\marginpar{95}
\begin{corollaire}
\label{IV.4.4}
Supongamos que~$A$ es un dominio \emph{local noetheriano}, de ideal
maximal~$\goth{m}=I$, cuerpo residual~$k=A/I$ y cuerpo de
fracciones~$K$. Sea~$M$ un módulo de tipo finito sobre~$A$. Entonces las
condiciones equivalentes anteriores~\textup{(i)} a~\textup{(v)} son
también equivalentes a
\ifthenelse{\boolean{orig}}{}
{\enlargethispage{.5cm}}%
\begin{enumerate}
\item[(vi)] $M\otimes_A K$ y~$M\otimes_A k$ son espacios vectoriales de
la misma dimensión (es decir, el rango de~$M$ sobre~$A$ es igual al número
mínimo de generadores del $A$-módulo~$M$).
\end{enumerate}
\end{corollaire}

La demostración es inmediata; se deja al lector la tarea de generalizar al
caso en que~$A$ solo se supone reducido: hay que exigir entonces que los
rangos de~$M$ en los ideales primos minimales de~$A$ sean iguales a la
dimensión del espacio vectorial~$M\otimes_A k$.
